\documentclass[a4paper]{exam}
\usepackage{amsmath,amssymb,titlesec}
\usepackage[a4paper,margin=22mm]{geometry}

\titleformat{\section}{\normalfont\Large\bfseries}{}{0pt}{}
\titleformat{\subsection}{\normalfont\large\bfseries}{}{0pt}{}
\titlespacing*{\section}{1em}{1.6\baselineskip}{0.6\baselineskip}
\titlespacing*{\subsection}{2em}{1.2\baselineskip}{0.4\baselineskip}

\setlength{\parindent}{0in}
\setlength{\headheight}{18pt}

\printanswers
\title{Week 12}
\author{Radhesh Goel}
\date{\today}

\begin{document}
\maketitle

\section*{Part 1 Theory}

\begin{solution}
  Write Burgers' equation in conservation form with flux
  \[
    f(u) = \frac{u^2}{2}
  \]
  The Rankine--Hugoniot condition gives the shock speed
  \[
    \begin{aligned}
      V_s & = \frac{f(u_L) - f(u_R)}{u_L - u_R}  \\
          & = \frac{u_L^2 - u_R^2}{2(u_L - u_R)} \\
          & = \frac{u_L + u_R}{2}
    \end{aligned}
  \]
  For \(u_L = 2\ \mathrm{m\,s^{-1}}\) and \(u_R = 1\ \mathrm{m\,s^{-1}}\),
  \[
    \boxed{V_s=1.5\ \mathrm{m\,s^{-1}}}
  \]
\end{solution}

\section*{Part 2 Programming}

\begin{solution}
  Both velocities are positive, so information travels from left to right and
  the upwind flux at node \(i\) comes from node \(i - 1\)
  \[
    \frac{u_i^{n+1} - u_i^n}{\Delta t} + \frac{(u_i^n)^2/2 - (u_{i-1}^n)^2/2}{\Delta x} = 0
  \]
  Solving for the new value gives the implemented update
  \[
    \boxed{
      u_i^{n+1}=u_i^n-
      \frac{\Delta t}{2\Delta x}
      \left[(u_i^n)^2-(u_{i-1}^n)^2\right]
    }
  \]
  Running the supplied problem with this update gives
  \[
    V_{s,\mathrm{conservative}} \approx 1.4993446\ \mathrm{m\,s^{-1}}
  \]
  This differs from the theoretical value by about \(0.0437\%\), whereas the
  original convective update gives
  \[
    V_{s,\mathrm{convective}} \approx 1.4145145\ \mathrm{m\,s^{-1}}
  \]
  which is about \(5.70\%\) low

  The conservative form is therefore much more accurate for the shock because it
  enforces the discrete flux balance and hence the Rankine--Hugoniot jump
  condition

  The convective and conservative differential forms are equivalent for smooth
  solutions, but the product-rule step connecting them is invalid across a
  discontinuity
\end{solution}

\section*{Practice Question 1}

\begin{solution}
  Taylor expansion about node \((i,j)\) gives
  \[
    \frac{u_{i+1,j} - 2u_{i,j} + u_{i-1,j}}{h^2} = u_{xx}
                                                   + \frac{h^2}{12}u_{xxxx}
                                                   + O(h^4)
  \]
  and similarly
  \[
    \frac{u_{i,j+1} - 2u_{i,j} + u_{i,j-1}}{h^2} = u_{yy}
                                                   + \frac{h^2}{12}u_{yyyy}
                                                   + O(h^4)
  \]
  The discrete PDE therefore represents
  \[
    u_{xx} + u_{yy} + \frac{h^2}{12}\left( u_{xxxx} + u_{yyyy} \right) + O(h^4) = 0
  \]
  Using the residual convention for local truncation error, its leading term is
  \[
    \boxed{
      \tau_{i,j}=\frac{h^2}{12}
      \left(u_{xxxx}+u_{yyyy}\right)
    }
  \]
  Differentiating \(u_{xx} + u_{yy} = 0\) twice also gives the equivalent form
  \[
    \boxed{\tau_{i,j}=-\frac{h^2}{6}u_{xxyy}}
  \]
  so the central-difference Laplacian is second-order accurate
\end{solution}

\section*{Practice Question 2}

\begin{solution}
  Let \(h_c\) denote the heat-transfer coefficient so it is not confused with
  the grid spacing used in Question 1

  A first-order backward difference uses the boundary node and the adjacent
  interior node
  \[
    \boxed{
      \frac{T_{3,3}-T_{3,2}}{\Delta y}
      =h_c\left(T_{3,3}-T_a\right)
    }
  \]
  Equivalently,
  \[
    \boxed{
      T_{3,3}
      =\frac{T_{3,2}-h_c\Delta y\,T_a}
      {1-h_c\Delta y}
    }
  \]
\end{solution}

\section*{Practice Question 3}

\begin{solution}
  The second-order backward approximation at the top boundary is
  \[
    \boxed{
      \frac{3T_{3,3}-4T_{3,2}+T_{3,1}}{2\Delta y}
      =h_c\left(T_{3,3}-T_a\right)
    }
  \]
  Solving for the boundary temperature gives
  \[
    \boxed{
      T_{3,3}
      =\frac{4T_{3,2}-T_{3,1}-2h_c\Delta y\,T_a}
      {3-2h_c\Delta y}
    }
  \]
\end{solution}

\end{document}
